\chapter{Introduction}
\label{ch:introduction}

\section{The Rise of Long-Context Large Language Models}
Large Language Models (LLMs) are no longer used only for short prompt completion. Current serving workloads often keep a document set, a code base, or a long dialogue history inside the same session. Model families such as LLaMA and Qwen now advertise context windows that would have looked unusual only a few years ago~\cite{llm_context, ring_attention}. The hardware consequence is easy to miss: a token admitted early in the session becomes part of the state that later tokens must revisit.

The checked traces pushed me away from treating this as only a matrix-throughput problem. For 2K and 8K contexts, the dense kernels are still visible enough that a faster arithmetic block matters. By 64K and 128K tokens, however, the repeated walk over old cache lines is the part that dominates my reading of the run. The serving system has to keep the KV cache resident, spill it when local HBM is short, and touch earlier tokens again for every decoded token. If that cache is sharded, the old state also has to agree across devices. That is why the later experiments focus on stored history and host-facing movement instead of only on a larger matrix unit.

\section{The Generative Memory Wall}
The serving path has two phases with very different memory behavior. In prefill, prompt tokens are processed together, and the matrix engines can exploit parallelism across the input sequence. Decode is the awkward phase for memory systems. Only one token is produced at a time, but that token still needs the \textit{Key-Value (KV) cache}: the query for the new token is scored against the keys and values saved from the entire prior history~\cite{vllm}.

The result is a state object whose size is tied directly to the usable context length. A 7B-class model with a $128{,}000$-token context can require several gigabytes of KV-cache storage for one request before batching or concurrent sessions are considered~\cite{vllm}. At that point, the problem is not simply how many multiply-accumulates the accelerator can issue; it is whether the memory hierarchy can hold the history and feed it back to the decode loop quickly enough.

Off-package memory expansion, including Compute Express Link (CXL) memory tiers, is a natural response when the cache no longer fits comfortably in local HBM~\cite{cxl_standard}. The catch is that capacity and movement are separate problems. A remote tier can hold the cache, but a conventional data path still pays to bring large portions of that cache back toward the host or accelerator during decode.

\section{Challenges in Existing Infrastructures}

\subsection{CXL and the Bandwidth Bottleneck}
Compute Express Link maps remote memory devices into the host's address space and makes pooled capacity easier to expose to software~\cite{cxl_pond}. For KV-cache serving, that abstraction is helpful but incomplete. The Root Complex and switch ports can become the busy part of the system, especially when the remote memory device is used as a passive store. In that host-aggregated path, the host or accelerator fetches KV tensors across the interconnect, performs attention-side work elsewhere, and repeats the transfer on the next token. With long contexts and many active requests, the repeated movement consumes shared bandwidth and creates queueing delay at the CXL switch ports~\cite{distserve}.

\subsection{Processing-in-Memory and Format Mismatches}
Processing-in-Memory (PIM) attacks the movement problem from the opposite direction: let memory-side logic consume the data before it leaves the package~\cite{samsung_pim}. For this thesis, the awkward part is numeric rather than conceptual. Serving stacks often keep KV caches in INT8 or INT4 forms to stretch effective capacity~\cite{llm_outliers, gptq}, while many PIM datapaths are designed around floating-point nonlinear units. If an endpoint first dequantizes KV data, then runs floating-point softmax, then reformats the result, the design spends too much of its locality advantage on conversion.

\section{Research Motivation and Objectives}
The design work in this report starts from a specific reading of the traces: for long-context decode, the limiting object is the KV cache, not the multiply-accumulate array. The useful question is where the cache is placed, how often it crosses an off-package link, and whether the nonlinear part of attention can run in the same compact numeric format used for storage.

Two design rules came out of this reading. First, a CXL endpoint should consume the KV cache in the format used by the serving stack; otherwise conversion becomes a second memory-side bottleneck. Second, a distributed memory pool should move only the softmax bookkeeping values that truly have to be global. The old K/V history is the expensive object, so a path that pulls it back to the host remains the baseline to beat.

\section{Claim Boundary}
I draw the boundary around the KV cache rather than around the whole Transformer. Model weights, dense matrix kernels, and the rest of the decoder stack remain on the accelerator side. The object allowed to move is the long-lived KV state, either into one CXL-attached HBM tier or across several CXL memory endpoints. This choice keeps the proposal independent of sparsity, attention approximation, prompt compression, and KV-cache compression. The narrower question is the one I can test with the artifact: if the cache has to remain available for future tokens anyway, how much host-facing traffic disappears when attention-side traversal runs beside that cache?

The evidence is separated for the same reason. There is no fabricated chip behind this report, so I do not present the result as one end-to-end silicon measurement. Single-endpoint timing comes from modified Ramulator runs; cluster behavior comes from the CXL-fabric event model; iNLU area and power come from synthesis proxy reports. The split is not pretty, but it keeps the assumptions visible. DRAM timing, switch queueing, and nonlinear datapath cost can be inspected separately instead of being hidden inside one speedup number.

\section{Thesis Contributions}
For the single-endpoint artifact, I built the datapath around one deliberately narrow job: let the CXL endpoint perform the attention work that repeatedly touches stored KV state, while the host and accelerator keep the rest of the generation pipeline. The endpoint receives a query descriptor, walks the HBM-resident K/V pages, and returns a compact attention result. The iNLU replaces the FP16 exponential path with range reduction, a quadratic fixed-point approximation, and shift scaling. The outlier front-end keeps ordinary activations on the dense INT8 lane and sends rare large values to a small INT32 buffer. In the 128K-token evaluation, this organization lowers read latency by more than 98.9\% and reduces normalized logic-die area by 17\% relative to the floating-point baseline.

DisaggKV is the pool-scale artifact. It starts from the case where a request's KV state is too large, or too shared, to treat as one endpoint's private region. Reused prefix pages are kept stable so later sessions can attach to them, while private continuation pages are spread enough to avoid turning one endpoint into the sink for a request's history. The reduction protocol then moves only two global softmax values across shards. In the modeled multi-tenant run, combined interconnect traffic falls by about 84\%, host-facing traffic by more than 90\%, the 50\,ms tail target holds until roughly 2.9K requests/s, and injected-fault recovery remains below one second.

\section{Thesis Organization}
The remaining chapters follow the order in which the design was built. Chapter~\ref{ch:background} records the attention, CXL, queueing, and PIM concepts used later. Chapter~\ref{ch:motivation} shows why the checked traces point to both capacity pressure and repeated movement. Chapters~\ref{ch:inlu_architecture} and~\ref{ch:disaggkv_system} describe the two artifacts: first one CXL-attached HBM endpoint, then a distributed CXL memory pool. Chapter~\ref{ch:methodology} explains the simulation inputs and provenance, Chapter~\ref{ch:evaluation} reports the measured trends, Chapter~\ref{ch:related_work} positions the work against nearby systems, and Chapter~\ref{ch:conclusion} closes with the remaining limitations.
